\section{Konstruktion von Anfangsschätzungen}
\subsection{Motivation}
Die Leistung von SQP-Verfahren hängt stark vom anfänglichen Parametervektor
ab. In Eisenbahnanwendungen liegen Anfangsdaten häufig als Messpolylinien,
bestehende Alignments oder teilweise strukturierte Entwurfsdaten vor.

\subsection{Eingabedarstellung}
\begin{definition}[Eingabepolylinie]
\[
P=(p_0,\dots,p_M),\quad p_i\in\R^2
\]
bezeichne eine gemessene oder importierte Polylinie.
\end{definition}

\subsection{Bogenlängenparametrisierung}
\begin{definition}
Definiere die kumulierte Bogenlänge
\[
s_0=0,\quad s_{i+1}=s_i+\|p_{i+1}-p_i\|.
\]
\end{definition}
Damit wird die Eingabe parametrisiert.

\subsection{Anfängliche Krümmungsschätzung}
\begin{definition}[Diskrete Krümmung]
Für drei aufeinanderfolgende Punkte definiere
\[
\kappa_i\approx\frac{\theta_{i+1}-\theta_i}{\Delta s_i},
\qquad\theta_i=\arg(p_{i+1}-p_i).
\]
\end{definition}
Dies liefert eine erste Krümmungsschätzung aus geometrischen Daten.

\subsection{Glättung}
Rohe Krümmungsschätzungen sind typischerweise verrauscht.
\begin{definition}[Glättungsoperator]
Wende einen Glättungsfilter an:
\[
\kappa_i^{(0)}=\sum_jw_{ij}\kappa_j.
\]
\end{definition}
Mögliche Verfahren sind gleitender Mittelwert, Gauß-Filter und Tiefpassfilter.

\subsection{Anfängliche Überhöhungsschätzung}
Überhöhung kann aus dem Gleichgewicht initialisiert werden:
\[
\phi_i^{(0)}=\arcsin\left(\frac{v^2\kappa_i^{(0)}}{g}\right).
\]
Alternativ kann sie nullgesetzt, aus Eingabedaten übernommen oder nach
Entwurfsregeln approximiert werden.

\subsection{Projektion auf das diskrete Gitter}
\begin{definition}
Interpoliere \(\kappa_i^{(0)},\phi_i^{(0)}\) auf das Optimierungsgitter
\[
s_0,\dots,s_N.
\]
\end{definition}

\subsection{Heuristische Segmentierung}
Das Eingabe-Alignment kann in Geraden, Kreisbögen und Übergangsbereiche
segmentiert werden.
\begin{definition}[Segmenterkennung]
Ein Segment wird anhand der Krümmung klassifiziert:
\[
\kappa_i\approx0\Rightarrow\text{Gerade},
\qquad
\kappa_i\approx\text{konstant}\Rightarrow\text{Kreisbogen}.
\]
\end{definition}
Diese Struktur verbessert die Anfangsschätzung und entspricht Ansätzen der
Alignment-Rekonstruktion nach Kufver
\cite{Kufver1997MathematicalDescription,Kufver2000RailwayAlignment}.

\subsection{Abschließende Anfangsschätzung}
\begin{definition}[Anfänglicher Parametervektor]
\[
\mathbf x_0=(\kappa_0^{(0)},\dots,\kappa_N^{(0)},
\phi_0^{(0)},\dots,\phi_N^{(0)}).
\]
\end{definition}

\subsection{Interpretation}
Die Anfangsschätzung verbindet gemessene Geometrie, Glättung und physische
Approximation.
\begin{summary}
Eine robuste Anfangsschätzung überführt geometrische Rohdaten in einen
strukturierten, optimierungsgeeigneten Parametervektor. Sie ist die wesentliche
Verbindung zwischen realen Daten und numerischen Verfahren.
\end{summary}
